\documentclass[12pt,a4paper]{article}

% ── Encodage et langue ──────────────────────────────────────────
\usepackage[utf8]{inputenc}
\usepackage[T1]{fontenc}
\usepackage[french]{babel}

% ── Mise en page ─────────────────────────────────────────────────
\usepackage[top=2.5cm, bottom=2.5cm, left=2.5cm, right=2.5cm]{geometry}
\usepackage{setspace}
\onehalfspacing

% ── Mathématiques ────────────────────────────────────────────────
\usepackage{amsmath, amssymb, amsthm}
\usepackage{mathtools}
\usepackage{bm}

% ── Tableaux ─────────────────────────────────────────────────────
\usepackage{booktabs}
\usepackage{tabularx}
\usepackage{array}
\usepackage{longtable}
\usepackage{multirow}
\usepackage{pdflscape}

% ── Typographie ──────────────────────────────────────────────────
\usepackage{microtype}
\usepackage{lmodern}
\usepackage{parskip}
\setlength{\parindent}{0pt}
\setlength{\parskip}{6pt}

% ── Liens ────────────────────────────────────────────────────────
\usepackage[colorlinks=true, linkcolor=black, citecolor=black,
            urlcolor=black]{hyperref}

% ── En-tête / pied de page ───────────────────────────────────────
\usepackage{fancyhdr}
\pagestyle{fancy}
\fancyhf{}
\fancyhead[L]{\small\textit{Méthodologie LOF --- Baseline de référence}}
\fancyhead[R]{\small\textit{MOUKOKO EKAMBI G.M.}}
\fancyfoot[C]{\thepage}
\renewcommand{\headrulewidth}{0.4pt}

% ── Sections ─────────────────────────────────────────────────────
\usepackage{titlesec}
\titleformat{\section}{\large\bfseries}{\thesection.}{0.5em}{}
\titleformat{\subsection}{\normalsize\bfseries}{\thesubsection}{0.5em}{}
\titleformat{\subsubsection}{\normalsize\itshape}{\thesubsubsection}{0.5em}{}

% ── Commandes utilitaires ────────────────────────────────────────
\newcommand{\nd}{--}

% ── Bibliographie ────────────────────────────────────────────────
\usepackage{natbib}

% ════════════════════════════════════════════════════════════════
\begin{document}

% ── Page de titre ────────────────────────────────────────────────
\begin{titlepage}
\centering
\vspace*{2cm}

{\LARGE\bfseries
DÉTECTION D'ANOMALIES DANS LES DONNÉES MONÉTAIRES BEAC-CEMAC\\[0.4em]
\large\normalfont
Modèle 1 : Local Outlier Factor comme baseline de référence ---
Justification, pipeline et comparaison avec BiVAT\\[0.4em]
}

\vspace{2cm}

{\large\bfseries MOUKOKO EKAMBI GEORGES MIGUEL}\\[0.5em]
{\normalsize\itshape
Stagiaire, Banque des États de l'Afrique Centrale (BEAC)\\
M2 Data Science \& IA --- ENSPD / Université de Douala, LIDIA Laboratory
}

\vspace{2cm}

{\normalsize Août 2026}

\vfill

{\small
Document rédigé dans le cadre du stage de recherche à la Banque des États
de l'Afrique Centrale (BEAC), M2 Data Science \& IA,
ENSPD / Université de Douala --- LIDIA Laboratory.
}
\end{titlepage}

% ── Résumé ───────────────────────────────────────────────────────
\begin{abstract}
Ce document présente la méthodologie complète du Modèle~1 d'une étude
comparative à deux niveaux portant sur la détection non supervisée de
variations atypiques dans les données monétaires des Autres Sociétés de
Dépôts (2SR) de la zone CEMAC. Le Modèle~1 est le Local Outlier Factor
(LOF), retenu comme baseline de référence parmi huit méthodes concurrentes
sur la base de critères empiriques documentés. Trois contributions structurent
ce document. Premièrement, une justification empirique du choix du LOF parmi
les méthodes de détection d'anomalies classiques, appuyée sur des benchmarks
publiés et cinq publications institutionnelles BIS/BCE/Banca d'Italia.
Deuxièmement, la documentation des cinq lacunes structurelles du LOF ---
ignorance des dépendances temporelles, dégénérescence en haute dimension,
absence de quantification d'incertitude, inexplicabilité par indicateur, et
instabilité des scores --- qui motivent le passage au Modèle~2 (BiVAT).
Troisièmement, un protocole de comparaison entre les deux modèles sans labels
ground truth, fondé sur des métriques d'évaluation non supervisée adaptées au
contexte de surveillance macroprudentielle BEAC-CEMAC.

\bigskip
\noindent\textbf{Mots-clés :} détection d'anomalies non supervisée, Local
Outlier Factor, densité locale, séries temporelles financières, BEAC, CEMAC,
pipeline de prétraitement, baseline comparative, BiVAT.
\end{abstract}

\tableofcontents
\newpage

% ════════════════════════════════════════════════════════════════
\section{Introduction}

\subsection{Cadre général de l'étude}

Ce document s'inscrit dans une étude comparative à deux modèles portant sur
la détection non supervisée de variations atypiques dans les données monétaires
de la Banque des États de l'Afrique Centrale (BEAC). Le Modèle~1, objet du
présent document, est le Local Outlier Factor (LOF), algorithme classique de
détection d'anomalies fondé sur la densité locale relative des observations.
Le Modèle~2, documenté séparément, est BiVAT (\textit{Bidirectional
VAE-Associative Transformer}), une architecture hybride profonde combinant
autoencodeur variationnel et Transformer bidirectionnel.

L'architecture de l'étude repose sur une logique de progression justifiée :
le LOF est choisi comme baseline parce qu'il est le meilleur compromis
disponible parmi les méthodes classiques non supervisées pour le contexte
BEAC (section~2), mais ses lacunes structurelles documentées (section~5)
motivent précisément la construction de BiVAT comme réponse architecturale
à chacune de ces lacunes. Le protocole de comparaison entre les deux modèles
(section~8) permet d'objectiver empiriquement ce gain sur les données BEAC
réelles en l'absence de labels d'anomalies certifiés.

\subsection{Le dataset BEAC-CEMAC}

Le dataset de référence regroupe les situations de l'actif et du passif des
Autres Sociétés de Dépôts (2SR) des six États membres de la CEMAC (Cameroun,
Congo, Gabon, Guinée équatoriale, RCA, Tchad), selon la nomenclature IFS du
FMI. Il se présente sous la forme de douze fichiers au format XLSX organisés
selon la nomenclature
\texttt{clean\_<pays>\_beac\_mapping\_2SR\_<Actif|Passif>.xlsx}.

Les caractéristiques structurelles de ce dataset conditionnent l'ensemble
des choix méthodologiques documentés dans ce travail.

\textit{Volume temporel.} Chaque fichier couvre environ 195 observations
mensuelles (décembre~2001 à mars~2026) pour le Cameroun, avec des couvertures
potentiellement plus courtes pour les pays à système bancaire moins développé
(RCA, Guinée équatoriale).

\textit{Dimensionnalité.} Environ 55 indicateurs financiers par fichier,
répartis en grandes catégories de bilan : Numéraire et Dépôts, Titres,
Crédits, Actions, Réserves Techniques d'Assurance, Autres Comptes à Recevoir
et Crédits Commerciaux. Les valeurs sont exprimées en milliards de FCFA.

\textit{Hétérogénéité inter-pays.} Les économies CEMAC présentent des
structures bancaires très hétérogènes : une même valeur absolue d'un agrégat
peut être parfaitement normale au Cameroun et extrême en RCA. Cette
hétérogénéité interdit toute comparaison normalisée directe et impose un
traitement séparé par pays.

\textit{Absence de labels certifiés.} Il n'existe pas d'étiquettes historiques
d'anomalies certifiées par un processus d'audit formalisé. La détection est
intrinsèquement non supervisée : les deux modèles produisent des scores
ordinaux dont la validation s'appuie sur des événements économiques documentés.

\textit{Ruptures structurelles documentées.} Le choc pétrolier 2014--2016,
la crise sécuritaire camerounaise à partir de 2017, le choc COVID-19 en
2020--2021 et les dynamiques post-pandémiques 2022--2024 constituent des
événements de référence pour la validation externe des détections.

% ════════════════════════════════════════════════════════════════
\section{Sélection de l'Algorithme : Pourquoi le LOF}

\subsection{Taxonomie des méthodes candidates}

La détection non supervisée d'anomalies dans des données multivariées repose
sur un ensemble de méthodes organisées en quatre familles selon le principe
de discrimination qu'elles mobilisent.

\textbf{Famille~1 --- Méthodes statistiques classiques.} Ces méthodes
supposent un modèle paramétrique de la distribution des données normales et
signalent les observations qui s'en écartent au-delà d'un seuil. Le
Z-score univarié, la distance de Mahalanobis multivariée et le seuillage
des résidus de modèles ARIMA ou VAR appartiennent à cette famille. Leur
limitation fondamentale est double : elles supposent la stationnarité de la
distribution des données normales, et elles supposent en général une
distribution paramétrique (gaussienne) que les données financières violent
systématiquement (queues épaisses, leptokurtose, régimes changeants).

\textbf{Famille~2 --- Méthodes à distance globale.} Ces méthodes évaluent
l'anomalie d'un point par sa distance aux autres observations dans un espace
global unique. La méthode kNN-distance (distance au k-ième plus proche voisin)
et le DB-outlier de Knorr \& Ng (1998) appartiennent à cette famille. Leur
limitation est l'utilisation d'une densité globale implicite qui ne s'adapte
pas aux variations locales de densité --- défaut critique pour des données
financières où les régimes pré et post-crise présentent des densités très
différentes dans l'espace des indicateurs.

\textbf{Famille~3 --- Méthodes à densité locale.} Ces méthodes évaluent
l'anomalie d'un point relativement à la densité de son voisinage immédiat,
permettant de détecter des outliers dans des données à densité variable. Le
LOF (Breunig et al.\ 2000), le COF (Tang et al.\ 2002), le LOCI (Papadimitriou
et al.\ 2003), le LoOP (Kriegel et al.\ 2009), l'INFLO (Jin et al.\ 2006)
et le LDOF (Zhang et al.\ 2009) appartiennent à cette famille.

\textbf{Famille~4 --- Méthodes d'ensemble et par isolement.} Ces méthodes
ne reposent pas sur la densité locale. L'Isolation Forest (Liu et al.\ 2008)
isole les anomalies par partitionnement aléatoire de l'espace. La One-Class
SVM (Schölkopf et al.\ 2001) apprend une frontière autour des données normales
dans un espace noyau. L'Elliptic Envelope suppose une distribution gaussienne
multivariée. DBSCAN en mode détection d'anomalies traite les points
non-assignés à un cluster comme outliers.

\subsection{Critères de sélection spécifiques au contexte BEAC}

Le choix de l'algorithme baseline est conditionné par cinq contraintes
structurelles du dataset BEAC qui éliminent certaines familles de méthodes
indépendamment de leurs performances générales.

\textit{Contrainte C1 --- Petit volume ($N = 195$).} L'Isolation Forest est
documenté comme sous-optimal sur des datasets de moins de 1\,000 observations
(Liu et al.\ 2008, ICDM) : son mécanisme d'isolement par partitionnement
aléatoire converge lentement sur petit volume. Les méthodes nécessitant
l'estimation d'une matrice de covariance de dimension $55 \times 55$
(Mahalanobis, Elliptic Envelope) sont numériquement instables quand
$N \approx 55 \times 3 = 165$ observations effectives après décomposition STL.

\textit{Contrainte C2 --- Dimension $D = 55$.} La malédiction de la
dimensionnalité affecte toutes les méthodes à base de distance. Zimek,
Schubert \& Kriegel (2012) identifient huit problèmes distincts pour les
méthodes à base de densité en haute dimension, dont la concentration des
distances qui rend le concept de voisinage non discriminant au-delà de
$D \approx 20$. Cette contrainte impose une réduction de dimension préalable
pour toutes les méthodes de la famille~2 et de la famille~3 (section~4.5).

\textit{Contrainte C3 --- Densité variable entre régimes.} Le dataset BEAC
présente des régimes macroéconomiques successifs avec des densités très
différentes dans l'espace des indicateurs : la période 2015--2019 (choc
pétrolier, ajustement) a une dynamique structurellement différente de la
période 2020--2023 (COVID-19, rebond). Les méthodes à distance globale
(Famille~2) et les méthodes supposant un cluster compact unique (One-Class
SVM) sont inadaptées à ce contexte. La densité locale relative du LOF est
structurellement adaptée aux données multi-régimes.

\textit{Contrainte C4 --- Absence de labels.} Toutes les méthodes supervisées
et semi-supervisées sont exclues d'emblée. Parmi les méthodes non supervisées,
celles qui nécessitent un ensemble de données normales certifié pour
l'entraînement (Deep SVDD, One-Class SVM avec noyau complexe) sont
difficiles à calibrer sans labels.

\textit{Contrainte C5 --- Exigence de reproductibilité institutionnelle.} Dans
le contexte BEAC/COBAC, la méthode doit être transparente, documentée et
reproductible. Les méthodes dont le comportement dépend de l'initialisation
aléatoire (Isolation Forest avec graines différentes) ou d'hyperparamètres
non documentés dans la littérature institutionnelle sont moins défendables
dans un rapport de surveillance macroprudentielle.

\subsection{Comparaison empirique des méthodes candidates}

\subsubsection{Résultats sur les datasets de référence}

Le tableau~\ref{tab:selection} présente les résultats empiriques comparatifs
des méthodes candidates sur les datasets de référence du domaine (SMD, MSL,
SMAP, SWaT, PSM) issus de Xu et al.\ (2022, ICLR, Table~1) et MEMTO
(arXiv:2312.02530, Table~1), ainsi que sur des données financières multivariées
réelles issues de arXiv:2209.11686 (Table~23--24).

\begin{longtable}{p{3.5cm}ccccccp{2.5cm}}
\caption{Comparaison empirique des méthodes candidates --- F1-score (\%) sur
datasets de référence et données financières.}
\label{tab:selection}\\
\toprule
\textbf{Méthode} & \textbf{SMD} & \textbf{MSL} & \textbf{SMAP} &
\textbf{SWaT} & \textbf{PSM} & \textbf{Fin.} & \textbf{Source} \\
\midrule
\endfirsthead
\toprule
\textbf{Méthode} & \textbf{SMD} & \textbf{MSL} & \textbf{SMAP} &
\textbf{SWaT} & \textbf{PSM} & \textbf{Fin.} & \textbf{Source} \\
\midrule
\endhead
\bottomrule\endlastfoot

Z-score & \nd & \nd & \nd & \nd & \nd & Très faible &
Univarié, exclu C2 \\
\midrule
Mahalanobis & \nd & \nd & \nd & \nd & \nd & Modérée &
Instable $N \approx D$, exclu C1 \\
\midrule
kNN-distance & \nd & \nd & \nd & \nd & \nd & \nd &
Densité globale, inférieur LOF \\
\midrule
DBSCAN & \nd & \nd & \nd & \nd & \nd & 26.46 &
arXiv:2209.11686 \\
\midrule
Isolation Forest & 53.64 & 66.45 & 55.53 & 47.02 & 83.48 & 17.82 &
Xu 2022 T.1 ; arXiv:2209.11686 \\
\midrule
OCSVM & 56.19 & 70.82 & 56.34 & 47.23 & 70.67 & 39.00 &
Xu 2022 T.1 ; arXiv:2209.11686 \\
\midrule
\textbf{LOF} & \textbf{46.68} & \textbf{61.18} & \textbf{57.60} &
\textbf{68.62} & \textbf{70.61} & \textbf{33.33} &
Xu 2022 T.1 ; arXiv:2209.11686 \\
\midrule
LoOP & \nd & \nd & \nd & \nd & \nd & \nd &
Variante LOF, voir \S\ref{subsec:loop} \\
\midrule
\textit{OmniAnomaly} & \textit{85.22} & \textit{87.67} & \textit{86.92} &
\textit{82.83} & \textit{80.83} & \nd &
\textit{Xu 2022 T.1 --- Modèle profond} \\
\midrule
\textit{BiVAT (cible)} & \textit{$\geq$92} & \textit{$\geq$93} &
\textit{$\geq$96} & \textit{$\geq$94} & \textit{$\geq$97} & \nd &
\textit{Modèle 2 de l'étude} \\
\end{longtable}

\textit{Note sur le protocole d'évaluation.} Les résultats sur SMD/MSL/SMAP/SWaT/PSM
utilisent le protocole \textit{point adjustment} (PA) : si au moins un point
d'un segment anormal est détecté, tout le segment est crédité. Ce protocole
tend à gonfler les scores par rapport à une évaluation strictement
point-à-point. Les résultats sur données financières (colonne Fin.) sont issus
d'un benchmark sur données de marché multivariées réelles sans point adjustment
(arXiv:2209.11686).

\subsubsection{Analyse comparative par famille}

\textbf{Méthodes statistiques classiques.} Le Z-score est structurellement
exclu par la contrainte C2 : il est univarié par construction et ne capture
pas les anomalies multivariées résultant de combinaisons anormales entre les
55 indicateurs IFS. Chandola, Banerjee \& Kumar (2009) formalisent cette
limitation : le Z-score détecte uniquement les anomalies ponctuelles
marginales, mais rate les anomalies contextuelles et collectives qui sont
précisément les plus pertinentes dans le contexte de surveillance
macroprudentielle. La distance de Mahalanobis est exclue par la contrainte
C1 : l'estimation stable d'une matrice de covariance $55 \times 55$ requiert
au minimum $55 \times 5 = 275$ observations, supérieur à $N = 195$. De plus,
comme méthode globale, elle suppose la stationnarité de la distribution normale,
violée par les changements de régime CEMAC. DOD (arXiv:2511.02199) confirme
que la distance de Mahalanobis ``degrades severely in high dimensions due to
the curse of dimensionality'' (Zimek et al.\ 2012).

\textbf{Méthodes à densité locale concurrentes du LOF.} Plusieurs variantes
du LOF sont disponibles dans la Famille~3. Le COF (Tang et al.\ 2002) utilise
un chemin de connectivité au lieu de la distance euclidienne, offrant un
avantage théorique sur les clusters non sphériques, mais son coût
computationnel $\mathcal{O}(n^2 k)$ supérieur à celui du LOF $\mathcal{O}(nk^2)$
est injustifié pour $N = 195$. Goldstein \& Uchida (2016) montrent que COF et
LOF ont des performances comparables sur données de dimension modérée. Le LOCI
(Papadimitriou et al.\ 2003) supprime le paramètre MinPts via un calcul
automatique, mais sa complexité $\mathcal{O}(n^2)$ en temps et en espace est
prohibitive pour $D = 55$. L'INFLO (Jin et al.\ 2006) intègre le voisinage
direct et inverse pour réduire les faux positifs aux frontières de clusters,
mais ne fournit pas de décomposition par indicateur. Le LDOF (Zhang et al.\
2009) est documenté comme plus stable que le LOF sur données éparses, mais
aucun benchmark sur données financières n'est disponible.

\textbf{Isolation Forest.} L'IF est compétitif sur les datasets industriels
(F1 = 83.48 sur PSM) mais inférieur au LOF sur les données financières
multivariées réelles (F1 = 17.82 vs 33.33, arXiv:2209.11686). Liu et al.\
(2008, ICDM original) documentent que l'IF est ``most accurate for data sets
larger than one thousand points'' : sur $N = 195$ (contrainte C1), son mécanisme
d'isolement par partitionnement aléatoire converge lentement et ses
performances se dégradent. Sur le dataset SWaT, le LOF dépasse l'IF de 21.60
points de F1 (68.62 vs 47.02, Xu et al.\ 2022), illustrant la supériorité
du LOF sur les données à densité hétérogène entre régimes. Goldstein \& Uchida
(2016) confirment que le LOF est généralement supérieur à l'IF sur les données
à densité variable.

\textbf{One-Class SVM.} L'OCSVM est compétitif sur MSL (F1 = 70.82) et sur
données financières (F1 = 39.00 vs 33.33 pour LOF sur arXiv:2209.11686), mais
il suppose que les données normales forment un cluster compact dans l'espace
noyau. Cette hypothèse est fondamentalement inadaptée aux données BEAC où la
distribution normale évolue entre régimes macroéconomiques distincts
(contrainte C3). L'OCSVM requiert également une calibration délicate du noyau
et du paramètre $\nu$ sans garantie de stabilité sur données non stationnaires.

\textbf{DBSCAN.} En mode détection d'anomalies, DBSCAN traite les points
non-assignés à un cluster comme outliers. Sur données financières multivariées,
sa F1 = 26.46 est inférieure à celle du LOF (33.33, arXiv:2209.11686). De
plus, DBSCAN produit une décision binaire (outlier/inlier) sans score continu,
ce qui interdit le classement ordinal des anomalies par sévérité --- exigence
opérationnelle fondamentale dans le contexte BEAC.

\subsection{Décision de sélection : le LOF comme baseline}

Le LOF est retenu comme algorithme baseline pour trois raisons complémentaires
et documentées.

\textit{Première raison --- Supériorité empirique dans le contexte
$N < 1\,000$ et données financières.} Sur données financières multivariées
réelles, le LOF obtient F1 = 33.33\%, supérieur à l'IF (17.82\%), DBSCAN
(26.46\%) et comparable à l'OCSVM (39.00\%, mais inadapté aux données
multi-régimes). Sur les datasets industriels à densité hétérogène (SWaT),
le LOF dépasse l'IF de 21.60 points de F1 (Xu et al.\ 2022). La combinaison
$N = 195$, $D = 55$ et densité variable entre régimes place le dataset BEAC
dans le domaine de compétence documenté du LOF.

\textit{Deuxième raison --- Légitimité institutionnelle documentée.} Le LOF
est explicitement référencé comme méthode de baseline dans au moins cinq
publications institutionnelles de surveillance bancaire : BIS IFC Bulletin~66
(2024) le cite comme référence pour la détection d'anomalies dans les données
bancaires ; Bank of Italy Occasional Paper~689 (Maddaloni et al., 2022) l'applique
aux données de reporting AnaCredit multi-sources ; BCE DG Statistics
(Accornero \& Boscariol, 2021) l'utilise sur données de reporting bancaire BCE ;
Banco de Portugal (Costa et al.\ 2021) l'applique au registre de crédit et
documente explicitement que les méthodes traditionnelles comme le LOF
``will (in general) perform poorly'' au-delà d'une certaine complexité, justifiant
le passage vers des méthodes plus sophistiquées ; BIS IFC Bulletin~57 (2022)
utilise des méthodes à densité (DBSCAN, Affinity Propagation) dans un contexte
opérationnel similaire.

\textit{Troisième raison --- Adéquation structurelle au problème de densité
variable.} La définition du LOF repose sur la densité locale relative : un
point est anomalie si sa densité locale est faible comparativement à celle
de ses voisins, indépendamment de la densité globale du dataset. Cette
propriété est précisément adaptée aux données BEAC où les régimes
macroéconomiques successifs créent des zones de densité très différentes dans
l'espace des 55 indicateurs. Les méthodes à densité globale (Z-score,
Mahalanobis, kNN-distance) mesurent l'écart à la distribution moyenne sur
toute la période, confondant ainsi les anomalies réelles avec les simples
changements de régime.

\subsection{LoOP : variante probabiliste non retenue}
\label{subsec:loop}

LoOP (\textit{Local Outlier Probabilities}, Kriegel et al.\ 2009,
DOI:~10.1145/1645953.1646195) est une variante directe du LOF qui transforme
les scores en probabilités d'anomalie dans $[0,1]$ via la fonction d'erreur
gaussienne appliquée au ratio PLOF/nPLOF. Son avantage principal est
l'interprétabilité directe du score comme probabilité, sans nécessiter de
seuillage préalable.

LoOP n'est pas retenu dans ce pipeline pour deux raisons précises. D'une part,
LoOP hérite de la même sensibilité à MinPts que le LOF standard : les
probabilités produites dépendent du choix du voisinage de la même manière que
les scores LOF bruts. D'autre part, l'interprétabilité probabiliste du score
est intégralement gérée dans ce pipeline par la Conformal Prediction post-hoc
(section~8.3), qui produit des intervalles de confiance à garantie formelle
$\Pr(\hat{y} \in C_\alpha) \geq 1 - \alpha$ --- supérieurs aux probabilités
LoOP qui supposent une distribution half-gaussienne des distances, hypothèse
non vérifiable sur données financières à régimes multiples. LoOP devient donc
redondant dans l'architecture pipeline retenue.
% ════════════════════════════════════════════════════════════════
\section{Cadre Théorique du LOF}

\subsection{\textit{k}-distance et voisinage}

Soit $\mathcal{D}$ l'ensemble des observations de l'espace résiduel normalisé
après décomposition STL et réduction par RPCA (sections~4.3 et~4.5). Pour un
entier $k \geq 1$ et un point $p \in \mathcal{D}$, la \textbf{\textit{k}-distance}
de $p$, notée $k\text{-dist}(p)$, est la distance $d(p, o)$ vers un objet
$o \in \mathcal{D}$ tel qu'au moins $k$ objets $o' \neq p$ vérifient
$d(p, o') \leq d(p, o)$ et au plus $k-1$ vérifient $d(p, o') < d(p, o)$
(Breunig et al.\ 2000, Déf.~3). Dans le contexte BEAC, $p$ représente un
mois $t$ dans l'espace des résidus réduit, et $d$ est la distance euclidienne
dans cet espace après réduction de dimension.

Le \textbf{voisinage de \textit{k}-distance} de $p$ est l'ensemble des points
situés à une distance inférieure ou égale à $k\text{-dist}(p)$ :
$$\mathcal{N}_k(p) = \bigl\{\, q \in \mathcal{D} \setminus \{p\}
\mid d(p, q) \leq k\text{-dist}(p) \,\bigr\}$$
(Breunig et al.\ 2000, Déf.~4). La cardinalité de $\mathcal{N}_k(p)$ peut
être supérieure à $k$ en cas d'équidistance de plusieurs points, phénomène
possible sur les données BEAC en raison de la faible variance de certains
indicateurs sur des sous-périodes de stabilité monétaire.

\subsection{Distance de joignabilité}

La \textbf{distance de joignabilité} (\textit{reachability distance}) de $p$
par rapport à $o$ (Breunig et al.\ 2000, Déf.~5) est :
$$\mathrm{reach\text{-}dist}_k(p, o) = \max \bigl\{
k\text{-dist}(o),\; d(p, o)
\bigr\}$$

Si $p$ est éloigné de $o$, la distance réelle $d(p, o)$ est retenue. Si $p$
est très proche de $o$ --- c'est-à-dire si $p$ se trouve à l'intérieur du
voisinage de $k$-distance de $o$ --- la distance est plafonnée à
$k\text{-dist}(o)$. Ce plafonnement réduit les fluctuations statistiques
locales induites par des configurations dégénérées où plusieurs points sont
presque coïncidents, situation possible sur les données BEAC pour des
indicateurs présentant de longues périodes de stagnation nominale.

\subsection{Densité locale de joignabilité}

La \textbf{densité locale de joignabilité} (\textit{local reachability density})
de $p$ (Breunig et al.\ 2000, Déf.~6) est l'inverse de la distance de
joignabilité moyenne vers les voisins de $p$ :
$$\mathrm{lrd}_k(p) = \left(
\frac{\displaystyle\sum_{o \in \mathcal{N}_k(p)}
\mathrm{reach\text{-}dist}_k(p, o)}
{\left|\mathcal{N}_k(p)\right|}
\right)^{-1}$$

Un $\mathrm{lrd}_k(p)$ élevé indique que $p$ est situé dans une région dense
de l'espace résiduel --- son comportement mensuel est similaire à celui de
plusieurs mois voisins, ce qui caractérise un comportement normal dans le
contexte BEAC. Inversement, un $\mathrm{lrd}_k(p)$ faible signale que $p$ est
isolé dans l'espace résiduel, première indication d'une déviation atypique.

\subsection{Facteur d'anomalie local}

Le \textbf{facteur d'anomalie local} de $p$ (Breunig et al.\ 2000, Déf.~7)
est le rapport moyen entre la densité locale de joignabilité des voisins de
$p$ et celle de $p$ lui-même :
$$\mathrm{LOF}_k(p) =
\frac{\displaystyle\sum_{o \in \mathcal{N}_k(p)}
\dfrac{\mathrm{lrd}_k(o)}{\mathrm{lrd}_k(p)}}
{\left|\mathcal{N}_k(p)\right|}$$

Un $\mathrm{LOF}_k(p) \approx 1$ caractérise un point normal : sa densité
locale est comparable à celle de ses voisins. Un $\mathrm{LOF}_k(p) \gg 1$
signale une anomalie locale : $p$ est situé dans une région beaucoup moins
dense que ses voisins immédiats. Appliqué aux données BEAC, un mois $t$
recevra un LOF élevé si ses résidus STL constituent une déviation simultanée
et atypique de plusieurs indicateurs du bilan par rapport à la dynamique
habituelle de ces indicateurs sur les mois voisins.

\subsection{Propriétés théoriques et bornes}

Breunig et al.\ (2000, Théorèmes~1 et~2) établissent des bornes générales
sur le facteur LOF en fonction des distances de joignabilité directes et
indirectes. Pour un point $p$ profondément intégré dans un cluster
$\mathcal{C}$ de densité uniforme, $\mathrm{LOF}_k(p)$ tend vers~1. Pour un
point $p$ isolé à une distance $d_{\min}$ du cluster le plus proche, le LOF
croît avec $d_{\min}$ et avec l'homogénéité du cluster environnant.

La propriété fondamentale du LOF pour le contexte BEAC est sa relativité
locale : deux points situés dans des régions de densités très différentes
(par exemple un mois de 2010 en régime de stabilité et un mois de 2021 en
régime post-COVID) peuvent recevoir le même score LOF s'ils sont également
anomaux relativement à leur contexte local respectif. Cette propriété est
précisément ce qui distingue le LOF des méthodes à distance globale et le
rend adapté aux données à régimes multiples.

\subsection{Stratégie de la plage de valeurs pour MinPts}

Le seul hyperparamètre du LOF est $\mathrm{MinPts}$, qui fixe le nombre de
voisins utilisé pour estimer la densité locale. Breunig et al.\ (2000,
Section~6.1) démontrent que le score LOF n'évolue pas de manière monotone
avec $\mathrm{MinPts}$ : pour un même point, augmenter $\mathrm{MinPts}$
peut alternativement augmenter, diminuer ou stabiliser son LOF selon la
structure locale du voisinage.

Pour pallier cette instabilité, Breunig et al.\ (2000, Section~6.2) proposent
d'utiliser une plage de valeurs $[\mathrm{MinPts}_\mathrm{LB},\,
\mathrm{MinPts}_\mathrm{UB}]$ et de retenir pour chaque point le score
maximum :
$$\mathrm{LOF}^*(p) = \max \bigl\{
\mathrm{LOF}_{\mathrm{MinPts}}(p) \mid
\mathrm{MinPts}_\mathrm{LB} \leq \mathrm{MinPts} \leq
\mathrm{MinPts}_\mathrm{UB}
\bigr\}$$

Le maximum met en évidence l'instance à laquelle le point est le plus isolé,
tandis que la moyenne diluerait ce caractère atypique. Pour le dataset BEAC
($N \approx 195$), la plage retenue est $[10, 30]$, correspondant à une borne
inférieure garantissant la stabilité statistique (règle pratique de Breunig
et al.) et à une borne supérieure de 30 mois (2,5 ans) cohérente avec la durée
des cycles macroéconomiques CEMAC. Pour les pays à volume réduit (RCA, Guinée
équatoriale), la borne supérieure est adaptée à
$\mathrm{MinPts}_\mathrm{UB} = \lfloor n/5 \rfloor$ afin de maintenir
un rapport $\mathrm{MinPts}/n$ cohérent avec les recommandations de Breunig
et al.

% ════════════════════════════════════════════════════════════════
\section{Pipeline d'Implémentation}

\subsection{Architecture multi-fichiers et mode d'exécution}

\subsubsection{Inventaire des sources}

Le dossier de données contient douze fichiers au format XLSX organisés selon
la nomenclature \texttt{clean\_<pays>\_beac\_mapping\_2SR\_<Actif|Passif>.xlsx},
couvrant les six États membres de la CEMAC (Cameroun, Congo, Gabon, Guinée
équatoriale, RCA, Tchad) sur les deux volets comptables Actif et Passif.
Chaque fichier est organisé en format large : chaque ligne représente un
indicateur financier identifié par son code IFS, et chaque colonne correspond
à une période mensuelle au format \texttt{AAAMyy}.

\subsubsection{Mode d'exécution séparé}

Trois architectures d'exécution sont envisageables : mode séparé (pipeline
complet exécuté indépendamment sur chacun des douze fichiers), mode pool
CEMAC (concaténation des douze fichiers avant exécution), et mode
hiérarchique (combinaison des deux niveaux). Le mode séparé est retenu pour
trois raisons structurelles.

Premièrement, la comparabilité inter-pays est structurellement inadéquate.
Les économies CEMAC présentent des structures bancaires très hétérogènes :
une même valeur absolue d'un agrégat peut être parfaitement normale au
Cameroun et extrême en RCA. Dans un pool CEMAC ($n \approx 195 \times 12 =
2\,340$), le LOF mesurerait l'écart à la distribution CEMAC composite plutôt
que l'anomalie propre à chaque pays, perdant toute signification économique.

Deuxièmement, la calibration de MinPts serait invalidée. Dans un pool de
2\,340 observations, le paramètre MinPts devrait être recalibré sur une base
douze fois plus grande, et son interprétation en termes de cycles de référence
annuels perdrait son ancrage économique.

Troisièmement, seul le mode séparé produit une interprétation économique
directe. Un score LOF élevé en avril~2020 au Cameroun doit être lu comme une
anomalie dans la dynamique bancaire camerounaise, et non comme une position
marginale dans un espace CEMAC composite. Le mode pool est conservé comme
variante d'exploration pour la détection d'anomalies systémiques simultanément
présentes dans plusieurs pays (section~\ref{sec:validation}).

\subsubsection{Cohérence comptable Actif--Passif}

L'identité fondamentale du bilan, $\mathrm{Total\;Actif} =
\mathrm{Total\;Passif}$, constitue un mécanisme de validation interne. Pour
chaque pays, les scores LOF obtenus sur le fichier Actif sont confrontés à
ceux obtenus sur le fichier Passif. Une anomalie concordante (signal sur les
deux volets au même mois) indique un événement économique réel affectant la
structure du bilan dans son ensemble. Une anomalie discordante (signal sur
un seul volet) suggère prioritairement une erreur de collecte ou un problème
de réconciliation comptable.

\subsection{Imputation des valeurs manquantes}

\subsubsection{Identification des types de lacunes}

Les fichiers sources présentent deux catégories de valeurs manquantes
aux mécanismes de génération distincts, ce qui conditionne le choix de
la méthode d'imputation.

Les \textit{lacunes internes courtes} (1 à 3 mois consécutifs) résultent
d'erreurs de calcul héritées des formules Excel sources ou de retards de
remontée. Elles sont raisonnablement supposées \textit{Missing At Random}
(MAR) : la probabilité d'une valeur manquante dépend de la période et de
l'indicateur (observables), non de la valeur non observée elle-même.

Les \textit{lacunes étendues internes} (plus de 3 mois consécutifs) résultent
de défaillances de collecte prolongées. L'hypothèse MAR reste raisonnable si
les indicateurs présentant des lacunes corrèlent avec d'autres indicateurs
renseignés.

Les \textit{lacunes structurelles en début ou fin de série} correspondent à
des périodes antérieures au démarrage effectif de la collecte d'un indicateur,
fréquentes en RCA et en Guinée équatoriale. Elles sont de nature
\textit{Missing Not At Random} (MNAR) : l'absence de valeur est liée à
l'indicateur lui-même et non à des facteurs observables. L'application de
MICE sur des données MNAR introduit un biais d'imputation sévère et non
contrôlé.

\subsubsection{Stratégie d'imputation par type de lacune}

\textit{Lacunes courtes (1 à 3 mois) :} interpolation temporelle linéaire.
Elle préserve la continuité locale de la série sans introduire de biais
structurel. Bezanson et al.\ (Nature Scientific Reports, 2026) confirment
que sur les séries temporelles à lacunes isolées, l'interpolation linéaire
``remained competitive across many settings'' comparativement aux méthodes
plus complexes. Son avantage supplémentaire dans le contexte de détection
d'anomalies est de ne pas exploiter les corrélations inter-séries pour
reconstruire la valeur manquante --- limitant ainsi le risque de lisser une
anomalie réelle présente simultanément dans plusieurs indicateurs.

\textit{Lacunes étendues (plus de 3 mois) :} algorithme MICE
(\textit{Multivariate Imputation by Chained Equations}, Van Buuren \&
Groothuis-Oudshoorn, 2011). MICE exploite les corrélations entre les 55
indicateurs du bilan pour inférer les valeurs absentes par régression
conditionnelle itérative sur chaque variable manquante. Cette approche est
adaptée au dataset BEAC car les indicateurs IFS présentent des corrélations
structurelles fortes (les sous-postes d'un agrégat covarient avec l'agrégat
parent). La validité de MICE repose sur l'hypothèse MAR, raisonnable pour
les lacunes étendues internes.

Un risque critique doit être signalé : toute méthode d'imputation peut masquer
des anomalies contenues dans les données manquantes. arXiv:2110.11780 documente
que les méthodes de gap filling ``tend to decrease the accuracy of the detection
results'' quand appliquées avant détection d'anomalies. Pour limiter ce risque,
les scores LOF obtenus sur les périodes comportant des valeurs imputées sont
signalés dans les résultats comme \textit{scores à fiabilité réduite}, distincts
des scores sur données complètes.

\textit{Lacunes structurelles MNAR :} troncature de la fenêtre temporelle
d'analyse au premier mois effectivement renseigné pour la majorité des
indicateurs. Les séries dont la part de valeurs imputées en début de période
dépasse 20\,\% de la fenêtre totale sont exclues de la matrice d'entrée du
LOF pour ce fichier.

\subsection{Normalisation robuste par série}

\subsubsection{Inadéquation du Z-score classique}

Les 55 indicateurs étant exprimés en milliards de FCFA mais présentant des
ordres de grandeur très hétérogènes, une normalisation individuelle par série
est indispensable avant tout calcul de distance euclidienne dans l'espace
multivarié. La standardisation classique $x^* = (x - \bar{x}) / \sigma$ est
inadaptée pour deux raisons complémentaires.

D'une part, la moyenne $\bar{x}$ et l'écart-type $\sigma$ sont très sensibles
aux valeurs extrêmes : un outlier prononcé gonfle simultanément $\bar{x}$ et
$\sigma$, réduisant mécaniquement sa propre valeur normalisée. L'anomalie
devient moins détectable après la normalisation censée préparer sa détection.
Leys et al.\ (2013) formalisent cette circularité sur données réelles et
montrent que le Z-score peut masquer des anomalies que la standardisation
robuste révèle.

D'autre part, dans le contexte BEAC, plusieurs indicateurs présentent des
ruptures structurelles documentées (choc COVID-19 2020--2021) qui constituent
précisément des anomalies à détecter. Le Z-score calculé sur l'ensemble de
la période 2001--2026 intégrerait ces ruptures dans l'estimation de $\sigma$,
biaisant la normalisation dans le sens contraire à l'objectif de détection.

\subsubsection{Standardisation retenue : médiane/MAD}

Pour chaque série $i$, la transformation appliquée est :
$$x_{i,t}^* = \frac{x_{i,t} - \tilde{x}_i}
                    {\kappa \cdot \mathrm{MAD}(x_i)}$$

où $\tilde{x}_i = \mathrm{median}_t(x_{i,t})$ est la médiane temporelle,
$$\mathrm{MAD}(x_i) =
\mathrm{median}_t\!\bigl(\,|x_{i,t} - \tilde{x}_i|\,\bigr)$$
est la déviation absolue médiane, et
$\kappa = 1/\Phi^{-1}(0{,}75) \approx 1{,}4826$ est le facteur de cohérence
rendant le MAD asymptotiquement équivalent à $\sigma$ sous distribution
gaussienne (Leys et al.\ 2013). La médiane possède un point de rupture de
50\,\% contre 0\,\% pour la moyenne, et le MAD hérite de cette propriété :
jusqu'à 50\,\% des observations peuvent être des anomalies sans biaiser
l'estimation de la dispersion de référence.

\subsection{Décomposition STL des séries temporelles}

\subsubsection{Inadéquation du LOF sur séries brutes}

Le LOF opère dans un espace géométrique statique : il évalue la densité locale
d'un point en fonction de ses distances aux voisins dans l'espace des
descripteurs, sans aucune modélisation de l'ordre temporel (lacune L1,
section~5). En présence de tendances macroéconomiques prononcées --- les
crédits aux sociétés non financières ont progressé de près de 600\,\% en
termes nominaux sur 2001--2026 --- la distance géométrique entre une
observation de 2010 et une observation de 2024 intègre l'effet de la dérive
temporelle, non l'anomalie conjoncturelle recherchée. De même, la saisonnalité
budgétaire récurrente (pics de fin d'année liés aux décaissements de l'État)
serait qualifiée d'anomalie par un LOF appliqué aux données brutes.

\subsubsection{Décomposition STL vs alternatives}

Trois méthodes de décomposition ont été évaluées pour isoler la composante
résiduelle stationnaire.

\textit{Différenciation première} : supprime les tendances stochastiques de
type I(1) mais ne gère pas la saisonnalité. Inadaptée aux données BEAC dont
la saisonnalité budgétaire est documentée et structurellement récurrente.

\textit{X-13 ARIMA-SEATS} : méthode standard des banques centrales pour les
statistiques officielles mensuelles/trimestrielles. Elle gère les effets
calendaires, les ajustements trading-day et produit des statistiques M1--M11
de qualité des résidus. Son inconvénient est qu'elle suppose une stationnarité
préalable et qu'elle est moins robuste aux outliers dans les séries d'entrée
--- les anomalies recherchées peuvent biaiser l'estimation des composantes.
Des recommandations techniques identifient X-13 comme adapté ``for monthly
and quarterly official statistics'' (revue metricgate.com/blogs/stl-vs-x13),
précisément le contexte BEAC.

\textit{STL (Seasonal and Trend decomposition using Loess, Cleveland et al.\
1990)} : décompose chaque série $x_{i,t}$ en trois composantes additives :
$$x_{i,t} = T_{i,t} + S_{i,t} + R_{i,t}$$
où $T_{i,t}$ est la tendance estimée par régression locale Loess, $S_{i,t}$
la composante saisonnière mensuelle (période $= 12$), et $R_{i,t}$ le résidu.

STL est retenu sur X-13 pour trois raisons spécifiques au contexte de
détection d'anomalies dans un pipeline ML. Premièrement, STL est robuste
aux outliers via un schéma de pondération interne (\textit{outer loop}) qui
atténue l'influence des points atypiques sur l'estimation de la tendance ---
évitant que les anomalies recherchées biaisent leur propre extraction.
Deuxièmement, sa composante saisonnière est \textit{time-varying}, adaptée
aux données dont la saisonnalité budgétaire peut évoluer avec les réformes
fiscales CEMAC. Troisièmement, STL s'intègre naturellement dans les pipelines
ML non supervisés selon les recommandations de la revue technique comparative.

Il est honnête de signaler qu'aucun benchmark quantitatif direct comparant
STL et X-13 sur des métriques de détection d'anomalies (F1, AUC) n'est
disponible dans la littérature publiée. Ce choix repose sur des arguments
qualitatifs documentés et sur la pratique établie des pipelines ML de
détection d'anomalies dans les séries temporelles. RobustSTL (AAAI 2019,
arXiv:1812.01767) constitue une amélioration du STL standard pour les séries
à saisonnalité variable et à bruit intense, mobilisable comme variante si
les résidus STL standards présentent une autocorrélation résiduelle
significative.

Seul le résidu $R_{i,t}$ est retenu comme entrée du LOF. Ce résidu représente
les fluctuations inexpliquées par la dynamique normale de la série et constitue
le support naturel de la détection d'anomalies.

\subsection{Réduction de la dimensionnalité par RPCA}

\subsubsection{Le fléau de la dimensionnalité pour le LOF}

Avec $D \approx 55$ indicateurs, l'espace résiduel dans lequel le LOF calcule
ses distances est de dimension~55. Zimek, Schubert \& Kriegel (2012) identifient
formellement huit problèmes distincts des méthodes à base de densité en haute
dimension. Le mécanisme central est la concentration des distances : en haute
dimension, les distances entre tous les points convergent vers la même valeur,
rendant le concept de voisinage non discriminant. CFOF (arXiv:1901.04992) prouve
formellement que kNN, aKNN et LOF ``concentrate according to Definition~3.5''
en haute dimension : les scores convergent vers une valeur constante, rendant
tout ranking impossible. Cette dégénérescence est une lacune structurelle du
LOF documentée en section~5.2.

Dans le contexte du dataset BEAC, ce problème est amplifié par la redondance
structurelle entre indicateurs : de nombreuses séries sont des composantes ou
des sous-totaux d'agrégats de niveau supérieur, induisant une forte
collinéarité qui concentre la variance dans un sous-espace de faible dimension.
Une réduction de dimension est donc indispensable avant le calcul du LOF, non
seulement pour contourner la malédiction de la dimensionnalité, mais aussi
pour éliminer la redondance informationnelle.

\subsubsection{RPCA vs ACP classique}

L'ACP classique est sensible aux outliers : ceux-ci exercent un effet de
levier sur les vecteurs propres, tirant les axes principaux vers eux et
biaisant la projection. Les anomalies se retrouvent alors projetées dans le
sous-espace jugé normal, réduisant leur visibilité dans l'espace de calcul
du LOF.

La RPCA (\textit{Robust Principal Component Analysis}, Candès et al.\ 2011)
contourne ce problème en décomposant la matrice $X \in \mathbb{R}^{n \times p}$
en une composante de bas rang $L$ (structure normale) et une composante creuse
$S$ (anomalies) :
$$X = L + S$$
par la minimisation convexe :
$$\min_{L,\,S} \; \|L\|_* + \lambda\|S\|_1
\quad \text{s.c.} \quad L + S = X$$
où $\|\cdot\|_*$ est la norme nucléaire et $\|\cdot\|_1$ la norme
\textit{entrywise} $\ell_1$, avec $\lambda = 1/\sqrt{\max(n,p)}$.

ROBPCA (PcaHubert, 2010) démontre que les méthodes robustes de type RPCA
``outperform several robust PCA tools'' et présentent ``the highest sensitivity''
pour la détection d'outliers. Sur données de réseau, arXiv:1801.01571 montre
que ``PCA analysis misses the vast majority of attack packets'' alors que RPCA
les détecte, au prix de faux négatifs significatifs pour l'ACP classique.

Seule la composante $L$ est ensuite projetée par ACP classique pour obtenir
les $k$ premières composantes principales, $k$ étant déterminé par le critère
du coude et validé par le pourcentage de variance expliquée cumulée (cible :
90\,\%). Il est honnête de signaler qu'aucun benchmark direct RPCA+LOF vs
ACP+LOF n'est disponible dans la littérature sur données financières ou
macroéconomiques multivariées. La justification de RPCA sur ACP repose sur
la décomposition $L + S$ qui isole mieux les anomalies ponctuelles ---
théoriquement justifiée (Candès et al.\ 2011) mais non quantifiée
empiriquement sur SMD/MSL/SMAP pour ce prétraitement spécifique.

arXiv:2603.18985 confirme que ``PCA models anomaly detection is limited to
linear relationships'' et qu'OmniAnomaly dépasse systématiquement les
approches PCA sur les benchmarks de référence, argument supplémentaire en
faveur de RPCA comme prétraitement robuste avant LOF, même si RPCA+LOF
reste structurellement inférieur aux méthodes profondes (lacune L2).

UMAP (\textit{Uniform Manifold Approximation and Projection}) est réservé à
la visualisation des scores LOF dans un espace bidimensionnel. Sa nature non
déterministe et la distorsion qu'il introduit dans les distances globales le
rendent inadapté comme espace de calcul du LOF dans un contexte opérationnel
où la reproductibilité est exigée (contrainte C5).

\subsection{Calcul du score LOF et formalisation du seuil}

\subsubsection{Score LOF agrégé sur la plage de MinPts}

Le pipeline produit, pour chaque mois $t$ de chaque fichier, un score agrégé :
$$\mathrm{LOF}^*(t) = \max \bigl\{
\mathrm{LOF}_{\mathrm{MinPts}}(t) \mid
10 \leq \mathrm{MinPts} \leq 30
\bigr\}$$

Ce score représente l'anomalie maximale observée sur l'ensemble de la plage
de voisinage, garantissant la robustesse aux fluctuations de MinPts
documentées en section~3.6.

\subsubsection{Seuil de détection par IQR}

Le seuil de détection $\tau$ est ancré dans la distribution empirique des
scores $\mathrm{LOF}^*$ via la règle de l'écart interquartile (Tukey, 1977) :
$$\tau = Q_3\!\bigl(\mathrm{LOF}^*\bigr)
+ 1{,}5 \times \mathrm{IQR}\!\bigl(\mathrm{LOF}^*\bigr)$$
où $Q_3$ est le troisième quartile et $\mathrm{IQR} = Q_3 - Q_1$.

Ce seuil est préféré au percentile fixe 95 car sa sensibilité ne dépend pas
de la proportion d'anomalies réelles dans le dataset --- inconnue a priori
dans le contexte non supervisé BEAC. En complément, le seuil au
95\textsuperscript{e} percentile est calculé comme variante de sensibilité.
Les observations dépassant les deux seuils simultanément constituent les
\textit{anomalies robustes}~; celles signalées par un seul seuil constituent
les \textit{anomalies marginales} requérant une expertise domaine.
% ════════════════════════════════════════════════════════════════
\section{Lacunes Structurelles du LOF}

Cette section cartographie les cinq lacunes structurelles du LOF qui
motivent la construction de BiVAT comme Modèle~2 de l'étude. Pour chaque
lacune, le mécanisme formel est exposé, les données empiriques documentées
sont présentées, et l'implication pour le contexte BEAC est explicitée.
La section~8 montrera comment BiVAT apporte une réponse architecturale
documentée à chacune de ces lacunes.

\subsection{Lacune L1 --- Ignorance des dépendances temporelles}

\subsubsection{Mécanisme formel}

Le LOF traite chaque observation comme un vecteur indépendant dans l'espace
des features. La définition du facteur d'anomalie local (section~3.4) ne fait
intervenir aucune notion d'ordre temporel, d'autocorrélation ou de rupture de
régime. Deux mois séparés par une décennie et présentant des résidus similaires
reçoivent exactement le même score LOF, indépendamment de leur position dans
le temps.

Cette limitation a une conséquence directe sur les types d'anomalies
détectables. Chandola, Banerjee \& Kumar (2009, ACM Computing Surveys, 41(3))
formalisent la taxonomie fondamentale des anomalies en séries temporelles :
les \textit{anomalies ponctuelles} (un seul instant $t$ aberrant dans l'espace
des features) sont détectables par le LOF ; les \textit{anomalies contextuelles}
(une valeur normale globalement mais anormale dans son contexte temporel
immédiat) sont ignorées car le LOF ne connaît pas ce contexte ; les
\textit{anomalies collectives} (une séquence d'observations normales
individuellement mais anormale en tant que séquence) sont également ignorées.
Or, dans le contexte de surveillance macroprudentielle BEAC, les anomalies
les plus pertinentes sont précisément contextuelles et collectives : un choc
de crédit en 2020 est anormal non parce que la valeur absolue des crédits est
hors norme, mais parce que sa trajectoire s'écarte brutalement de la dynamique
attendue.

\subsubsection{Données empiriques}

Le benchmark MEMTO (Xu et al.\ 2023, arXiv:2312.02530, Table~1) fournit la
comparaison la plus complète disponible entre le LOF et les méthodes à
modélisation temporelle sur les cinq datasets de référence.

\begin{table}[ht]
\centering
\caption{LOF vs méthodes temporelles profondes --- F1-score (\%).
Source : Xu et al.\ 2022 (ICLR, Table~1) et MEMTO (arXiv:2312.02530,
Table~1). Protocole point adjustment.}
\label{tab:l1-benchmark}
\small
\begin{tabular}{lcccccr}
\toprule
\textbf{Méthode} & \textbf{SMD} & \textbf{MSL} & \textbf{SMAP} &
\textbf{SWaT} & \textbf{PSM} & \textbf{Moyenne} \\
\midrule
LOF & 46.68 & 61.18 & 57.60 & 68.62 & 70.61 & 60.94 \\
OmniAnomaly (LSTM-VAE) & 85.22 & 87.67 & 86.92 & 82.83 & 80.83 & 84.69 \\
Anomaly Transformer & 92.33 & 93.59 & 96.69 & 94.07 & 97.89 & 94.91 \\
\midrule
\textbf{Écart LOF vs AT} &
$-$45.65 & $-$32.41 & $-$39.09 & $-$25.45 & $-$27.28 &
\textbf{$-$33.98} \\
\bottomrule
\end{tabular}
\end{table}

L'écart moyen entre le LOF et l'Anomaly Transformer est de $-33{,}98$ points
de F1. Ce n'est pas un écart de réglage fin --- c'est un écart de nature :
les méthodes temporelles détectent des anomalies que le LOF rate
structurellement parce qu'elles modélisent les dépendances séquentielles.
arXiv:2208.09240 confirme ce mécanisme : les méthodes non-temporelles
``perform the worst because they are poor at capturing temporal information.''
DAGMM (modèle génératif sans dépendance temporelle explicite) obtient
F1 = 57.30 sur SMD, très proche du LOF (46.68), confirmant que l'absence
de modélisation temporelle est le facteur limitant commun.

\subsubsection{Implication pour le contexte BEAC}

Sur les données BEAC, les anomalies les plus économiquement significatives
sont de nature contextuelle et collective : le choc COVID-19 n'est pas un
outlier ponctuel dans l'espace des 55 indicateurs, mais une rupture de
trajectoire sur plusieurs mois consécutifs. Le LOF appliqué aux résidus STL
peut détecter des déviations ponctuelles importantes (erreurs de saisie,
chocs exogènes brutaux) mais sous-détectera systématiquement les dérives
progressives et les anomalies collectives de type choc macroéconomique étendu.

\subsection{Lacune L2 --- Dégénérescence en haute dimension}

\subsubsection{Mécanisme formel}

Zimek, Schubert \& Kriegel (2012) identifient le phénomène de concentration
des distances comme le problème central des méthodes à base de densité en
haute dimension : lorsque $D$ croît, la distribution des distances entre
tous les points converge vers une valeur constante, rendant le classement
par proximité non discriminant. CFOF (arXiv:1901.04992) prouve formellement
que kNN, aKNN et LOF ``concentrate according to Definition~3.5'' en haute
dimension --- les scores LOF convergent vers une valeur constante, rendant
tout ranking d'anomalies impossible. Cette concentration est une propriété
théorique universelle des méthodes à base de distance euclidienne en espace
de grande dimension.

\subsubsection{Données empiriques}

Sur données financières multivariées réelles (arXiv:2209.11686, Table~23--24),
le LOF obtient F1 = 33.33\% contre 56.35\% pour PCA + réseau de neurones,
illustrant que même après réduction de dimension, la chaîne LOF reste inférieure
aux méthodes exploitant des représentations non-linéaires. arXiv:2603.18985
confirme que ``PCA models anomaly detection is limited to linear relationships''
et qu'OmniAnomaly dépasse systématiquement les approches PCA sur les benchmarks
de référence. La réduction RPCA atténue le problème mais ne le résout pas :
la projection dans un sous-espace linéaire de bas rang ne capture pas les
structures non-linéaires des anomalies financières complexes.

\subsubsection{Implication pour le contexte BEAC}

Avec $D = 55$ indicateurs, le dataset BEAC se situe au-delà du seuil de
dégénérescence documenté par Zimek et al.\ (2012) pour les méthodes à densité
locale ($D \approx 20$). La réduction RPCA ramène la dimension effective
à $k < 20$ composantes (cible de 90\,\% de variance expliquée), atténuant
partiellement ce problème. Mais les relations non-linéaires entre indicateurs
IFS --- notamment les interactions entre crédits, dépôts et liquidité dans
les phases de stress bancaire --- ne sont pas capturées par une projection
linéaire, quelle qu'elle soit.

\subsection{Lacune L3 --- Absence de quantification d'incertitude}

\subsubsection{Mécanisme formel}

Le LOF produit un score scalaire déterministe dans $[0, +\infty)$ sans
intervalle de confiance ni probabilité associée. Ce score n'est pas
directement interprétable comme probabilité d'anomalie : deux mois recevant
des scores de 1.8 et 2.1 peuvent être également anormaux ou très différemment
anormaux selon la distribution empirique des scores sur le dataset, sans que
le score lui-même permette de le déterminer.

\subsubsection{Données empiriques et variante LoOP}

La variante LoOP (Kriegel et al.\ 2009) transforme les scores LOF en
probabilités dans $[0,1]$ via la fonction d'erreur gaussienne appliquée
au ratio PLOF/nPLOF. Son avantage est l'interprétabilité directe du score,
mais elle suppose une distribution half-gaussienne des distances --- hypothèse
non vérifiable sur données financières à régimes multiples. De plus, LoOP
hérite de la même sensibilité à MinPts que le LOF standard. BIS FSI Paper
(2024) souligne que ``the lack of explainability of certain AI models can make
it challenging for regulators to ascertain financial institutions' compliance''
--- un score sans incertitude rend également impossible le contrôle du taux
de fausse alarme, exigence fondamentale de la surveillance prudentielle.

\subsubsection{Implication pour le contexte BEAC}

Un superviseur prudentiel BEAC a besoin non seulement d'un score scalaire
indiquant qu'une observation est anormale, mais d'une mesure de la certitude
de cette décision pour calibrer la priorité des investigations. Un score LOF
de 2.5 en avril 2020 ne dit pas si cette anomalie est certaine (haute confiance)
ou marginale (à surveiller). La Conformal Prediction post-hoc, appliquée dans
BiVAT, apporte cette garantie formelle $\Pr(\hat{y} \in C_\alpha) \geq 1 -
\alpha$ que le LOF standard ne peut pas fournir sans post-traitement.

\subsection{Lacune L4 --- Absence d'explicabilité par indicateur}

\subsubsection{Mécanisme formel}

Le LOF produit un score global par observation. Il ne décompose pas la
contribution de chaque indicateur IFS à l'anomalie détectée. Pour un mois
$t$ recevant un score $\mathrm{LOF}^*(t) = 3.2$, il est impossible de
déterminer si cette anomalie est portée principalement par les crédits aux
entreprises, par les dépôts transférables, ou par une combinaison de plusieurs
indicateurs --- sans recourir à un post-traitement externe au modèle lui-même.

\subsubsection{Extensions documentées et leurs limites}

Les extensions du LOF proposant une explicabilité par feature sont rares et
insuffisantes. La méthode Feature Bagging (Lazarevic \& Kumar, 2005, KDD)
agrège des scores LOF calculés sur des sous-espaces aléatoires de features,
permettant d'identifier les sous-espaces les plus discriminants. Cependant,
elle ne fournit pas une attribution quantitative par feature individuelle
directement interprétable. L'INFLO (Jin et al.\ 2006) et le LDOF (Zhang et al.\
2009) améliorent la stabilité et la robustesse du LOF mais n'ajoutent pas de
décomposition par indicateur. La seule solution d'explicabilité feature-level
applicable au LOF est SHAP post-hoc appliqué non pas au modèle lui-même
(le LOF n'est pas différentiable) mais aux features d'entrée --- approche
approximative qui ne satisfait pas les exigences d'auditabilité de SHAPAttenVAE
(Zhang et al.\ 2025) où l'explicabilité est intégrée à l'architecture.

\subsubsection{Implication réglementaire pour le contexte BEAC}

L'ECB Guide to Internal Models (2025, ML chapter) et le BIS FSI Paper (2024)
exigent une attribution causale par feature reproductible et auditable pour
les modèles de surveillance prudentielle. SHAP est la seule méthode
explicitement citée dans les processus ICAAP/CCAR à ce titre (SHARC,
arXiv:2607.05484). Dans le contexte BEAC/COBAC, l'impossibilité de produire
une explication du type ``l'anomalie d'avril 2020 est attribuée à 67\% à la
série Crédits\_ENF et 22\% à Dépôts\_transf'' constitue une limitation
opérationnelle majeure du LOF qui empêche son déploiement comme outil de
surveillance autonome sans expertise humaine complémentaire.

\subsection{Lacune L5 --- Instabilité des scores selon MinPts}

\subsubsection{Mécanisme formel}

Breunig et al.\ (2000, Section~6.1) démontrent dès le papier original la
non-monotonicité du score LOF en fonction de MinPts : pour un même point,
augmenter MinPts peut alternativement augmenter, diminuer ou stabiliser son
LOF selon la structure locale du voisinage. Il n'existe pas de valeur
universellement optimale de MinPts, et le score d'un point peut changer
radicalement selon le voisinage considéré.

\subsubsection{Données empiriques}

Goldstein \& Uchida (2016, PLoS ONE, 11(4): e0152173) conduisent un benchmark
systématique de 19 algorithmes de détection d'anomalies sur données
multivariées avec variation des hyperparamètres. Leur conclusion est que le
LOF est parmi les algorithmes les plus sensibles aux hyperparamètres sur les
datasets à densité variable. En comparaison, Liu et al.\ (2008, ICDM) montrent
que l'Isolation Forest ``converges very quickly at small $\psi$'' et que
``AUC is near optimal when $\psi = 128$'' --- l'IF converge avec peu d'arbres
et est stable dans une large plage de paramètres, avantage documenté sur le LOF.

\subsubsection{Implication pour le contexte BEAC}

L'instabilité des scores LOF selon MinPts est partiellement atténuée par la
stratégie de la plage de valeurs (section~3.6), qui calcule le maximum sur
$\mathrm{MinPts} \in [10, 30]$. Cette stratégie réduit la sensibilité mais
ne l'élimine pas : le score $\mathrm{LOF}^*(t)$ dépend toujours des bornes
de la plage retenue. Dans un contexte de surveillance institutionnelle où la
reproductibilité est exigée (contrainte C5), cette dépendance au choix de la
plage doit être documentée explicitement et justifiée dans tout rapport de
résultats.

% ════════════════════════════════════════════════════════════════
\section{Protocole de Validation}

\subsection{Contexte non supervisé}

La détection d'anomalies dans le dataset BEAC est intrinsèquement non
supervisée : il n'existe pas de labels historiques certifiés. Breunig et al.\
(2000, Section~7) soulignent cette dimension exploratoire du LOF : il s'agit
d'un outil de scoring ordinal dont la valeur réside dans le classement relatif
des observations. L'évaluation des détections repose donc sur cinq niveaux
complémentaires de validation externe.

\subsection{Protocole en cinq niveaux}

\textbf{Niveau 1 --- Analyse de sensibilité des scores LOF.}
Le profil $\mathrm{LOF}_{\mathrm{MinPts}}(t)$ est tracé pour
$\mathrm{MinPts} \in [10, 30]$, reproduisant la démarche de la Figure~8
de Breunig et al.\ (2000) sur les données BEAC. Les mois maintenant un LOF
élevé sur toute la plage sont des anomalies robustes au sens de la
section~4.6. Les mois anomaux seulement pour certaines valeurs de MinPts
sont signalés comme anomalies sensibles requérant une expertise domaine.

\textbf{Niveau 2 --- Confrontation aux ruptures macroéconomiques documentées.}
Les anomalies détectées ($\mathrm{LOF}^* > \tau$) sont confrontées aux
événements documentés dans la zone CEMAC : choc pétrolier 2014--2016, crise
sécuritaire camerounaise à partir de 2017, choc COVID-19 2020--2021, dynamiques
post-pandémiques 2022--2024. Une concordance entre les anomalies détectées et
ces épisodes constitue une validation externe qualitative. Cette approche est
cohérente avec la méthode utilisée par BIS Bulletin~84 (2024), qui valide les
détections d'autoencoders sur systèmes de paiement par concordance avec
``real-life anomalies including operational disruptions among important domestic
banks''.

\textbf{Niveau 3 --- Comparaison avec une baseline de distance globale.}
Le LOF est comparé à la distance de Mahalanobis multivariée appliquée aux
mêmes résidus STL normalisés, après vérification de la stabilité numérique
de l'estimation de la matrice de covariance. Ce niveau met en évidence les
anomalies locales que le LOF détecte et que la méthode globale manque ---
notamment les déviations atypiques dans des régimes à densité variable
(post-2020 vs pré-2015). La métrique de comparaison est la concordance entre
les top-20 anomalies classées par chaque méthode.

\textbf{Niveau 4 --- Validation croisée inter-pays.}
Les douze jeux de scores $\mathrm{LOF}^*$ sont superposés sur l'axe temporel.
Un mois $t$ est qualifié d'\textbf{anomalie systémique CEMAC} s'il dépasse
$\tau$ dans au moins $m \geq 3$ pays simultanément. Les anomalies systémiques
devraient coïncider avec les chocs documentés (crise pétrolière, COVID-19).
Une anomalie détectée dans un seul pays sur une période non documentée est
plus vraisemblablement liée à un problème de collecte local qu'à un événement
économique réel.

\textbf{Niveau 5 --- Cohérence Actif--Passif.}
Les scores $\mathrm{LOF}^*$ des fichiers Actif et Passif sont comparés mois
par mois pour chaque pays. Les anomalies concordantes sur les deux volets
sont interprétées comme des événements économiques réels~; les anomalies
discordantes sont signalées comme suspectes et soumises à vérification
comptable.

% ════════════════════════════════════════════════════════════════
\section{Positionnement Institutionnel du LOF}

Cette section documente la légitimité du LOF comme baseline de référence
dans le contexte de surveillance bancaire et financière institutionnelle,
en s'appuyant sur cinq publications directement pertinentes.

\textbf{BIS IFC Bulletin~66 (2024).} Ce bulletin cite explicitement le LOF
(``LOF: Identifying Density-Based Local Outliers'', Breunig et al.\ 2000)
comme méthode de référence dans le contexte de détection d'anomalies dans les
données bancaires, aux côtés de l'Isolation Forest et d'ECOD.
\url{https://www.bis.org/ifc/publ/ifcb66\_08.pdf}

\textbf{Bank of Italy Occasional Paper~689 (2022).} Maddaloni et al.\ (SSRN:
4109500) appliquent des algorithmes ML incluant des méthodes de la famille LOF
pour détecter des anomalies dans les données de reporting bancaire granulaires
(AnaCredit vs BSI vs FinRep). Le papier rapporte que ``performance of the
stacking technique, in terms of F1-score, is higher than in each algorithm
alone'', validant l'approche comparative multi-méthodes adoptée dans ce travail.

\textbf{BCE / ECB DG Statistics (2021).} Accornero \& Boscariol (BIS IFC
Bulletin~57, Section~5) appliquent des méthodes ML à la détection d'anomalies
dans les données de reporting bancaire à la BCE, dans un cadre non supervisé
comparable au contexte BEAC.
\url{https://www.bis.org/ifc/publ/ifcb57\_05.pdf}

\textbf{Banco de Portugal (2021).} Costa, Fonseca \& Maurício (BIS IFC
Bulletin~57, Section~29) appliquent des méthodes de détection d'anomalies
au registre de crédit portugais et documentent explicitement que ``more
traditional outlier detection methods will (in general) perform poorly,
especially from the computational point of view'' au-delà d'une certaine
complexité des données. Cette citation fournit la justification institutionnelle
du passage vers des méthodes plus sophistiquées comme BiVAT.
\url{https://www.bis.org/ifc/publ/ifcb57\_29.pdf}

\textbf{BIS IFC Bulletin~57 (2022).} Ce bulletin recense les applications ML
dans les banques centrales et décrit l'utilisation opérationnelle de méthodes
à densité (Affinity Propagation + DBSCAN) pour la détection d'anomalies dans
les données reportées par les institutions financières --- même famille
algorithmique que le LOF.
\url{https://www.bis.org/ifc/publ/ifcb57\_01\_rh.pdf}

% ════════════════════════════════════════════════════════════════
\section{Comparaison LOF vs BiVAT}

\subsection{Tableau comparatif structurel}

Le tableau~\ref{tab:lof-bivat} présente la comparaison structurelle entre
le LOF (Modèle~1) et BiVAT (Modèle~2) sur les cinq lacunes identifiées en
section~5. Pour chaque lacune, la colonne BiVAT indique la réponse
architecturale apportée et sa justification empirique documentée dans le
document d'état de l'art BiVAT.

\begin{longtable}{p{3cm}p{5.5cm}p{6cm}}
\caption{Comparaison structurelle LOF vs BiVAT sur les cinq lacunes.}
\label{tab:lof-bivat}\\
\toprule
\textbf{Lacune} & \textbf{LOF (Modèle 1)} & \textbf{BiVAT (Modèle 2)} \\
\midrule
\endfirsthead
\toprule
\textbf{Lacune} & \textbf{LOF (Modèle 1)} & \textbf{BiVAT (Modèle 2)} \\
\midrule
\endhead
\bottomrule\endlastfoot

L1 --- Dépendances temporelles &
Aucune modélisation temporelle. Anomalies contextuelles et collectives
non détectables. Écart moyen $-33{,}98$ pts F1 vs Anomaly Transformer
(Xu et al.\ 2022). &
Encodeur hybride Transformer bidirectionnel + BiLSTM. Modélisation des
dépendances longue portée et locales. Score collectif sur fenêtre glissante.
F1 moyen $\geq 94$\% sur benchmarks de référence. \\
\midrule

L2 --- Haute dimension &
Dégénérescence documentée au-delà de $D \approx 20$ (Zimek et al.\ 2012).
RPCA atténue mais ne résout pas le problème (relations non-linéaires
non capturées). &
Espace latent VAE capturant les relations non-linéaires entre les 55
indicateurs. Encoder contextuel bidirectionnel dans l'espace latent.
Supérieur aux approches PCA+méthode classique (arXiv:2603.18985). \\
\midrule

L3 --- Incertitude &
Score déterministe sans intervalle de confiance. LoOP exclu (hypothèse
half-gaussienne invérifiable, sensibilité MinPts identique). &
Conformal Prediction post-hoc : garantie formelle
$\Pr(\hat{y} \in C_\alpha) \geq 1-\alpha$. SSBC pour petit calibration set
($N=195$). CP pondérée pour non-stationnarité. \\
\midrule

L4 --- Explicabilité &
Score global sans décomposition par indicateur IFS. Feature Bagging
insuffisant. SHAP post-hoc approximatif (LOF non différentiable). &
SHAP DeepSHAP inter-couches intégré à l'architecture. Perte de cohérence
$\mathcal{L}_{\mathrm{expl}}$. Attribution par indicateur IFS directement
exploitable. Stabilité SHAP = 1.00 (arXiv:2012.04218). \\
\midrule

L5 --- Stabilité paramétrique &
Non-monotonicité documentée (Breunig et al.\ 2000). Parmi les algorithmes
les plus sensibles aux hyperparamètres (Goldstein \& Uchida, 2016). Stratégie
de plage atténue sans éliminer. &
Hyperparamètres appris par optimisation end-to-end. KL annealing pour
stabiliser l'entraînement variationnel. Association Discrepancy : $+18{,}82$
pts F1 sans surcoût computationnel (ablation Xu et al.\ 2022, Table~3). \\
\end{longtable}

\subsection{Protocole de comparaison empirique sans labels}

En l'absence de labels ground truth, la comparaison empirique entre LOF et
BiVAT sur les données BEAC repose sur quatre métriques d'évaluation non
supervisée complémentaires.

\textbf{Métrique M1 --- Precision@k sur événements documentés.}
Pour les $k$ mois recevant les scores les plus élevés (LOF$^*$ ou score BiVAT),
on calcule la proportion correspondant à des événements macroéconomiques
documentés dans la zone CEMAC (liste de référence : choc pétrolier
2014--2016, crise camerounaise 2017--2019, COVID-19 mars--décembre 2020,
rebond 2022). Cette métrique opérationnalise la validation par concordance
utilisée par BIS Bulletin~84 (2024) et BIS IFC Bulletin~66 (2024). Une
Precision@k plus élevée indique que le modèle classe les mois économiquement
documentés plus haut dans son ranking d'anomalies.

\textbf{Métrique M2 --- Rang moyen des anomalies connues.}
Pour chaque événement documenté (COVID-19 mars 2020, nadir de la crise
pétrolière décembre 2015, pic de la crise de liquidité 2016), on relève le
percentile du score d'anomalie à cette date dans la distribution empirique
des scores sur la période complète. Un rang au 95\textsuperscript{e} percentile
ou supérieur valide que le modèle détecte l'événement comme anomalie robuste.

\textbf{Métrique M3 --- Cohérence inter-modèles sur les top-k alertes.}
Le taux de concordance entre les top-20 alertes LOF et les top-20 alertes
BiVAT est calculé pour chaque fichier pays/volet. Une concordance élevée
($\geq 70\,\%$) indique que les deux modèles s'accordent sur les anomalies
les plus sévères et renforce la robustesse des détections. Une faible
concordance indique que les deux modèles détectent des types d'anomalies
différents (ponctuelles pour LOF, contextuelles/collectives pour BiVAT),
information utile pour la surveillance.

\textbf{Métrique M4 --- Concordance inter-pays sur anomalies systémiques.}
Pour les anomalies classées systémiques CEMAC (section~6, Niveau~4), on
compare le taux de concordance entre LOF et BiVAT. Les anomalies systémiques
confirmées par les deux modèles constituent les alertes de niveau 1 (à
signaler immédiatement)~; celles détectées par un seul modèle constituent
les alertes de niveau 2 (à analyser).

\subsection{Métriques communes et présentation des résultats}

Les deux modèles produisent leurs résultats dans un format standardisé commun
permettant la comparaison directe. Pour chaque fichier pays/volet et chaque
mois $t$, les résultats incluent le score d'anomalie normalisé (LOF$^*$
ou score BiVAT discriminant), le niveau de confiance associé (IQR-rule pour
LOF, intervalle CP pour BiVAT), la classification de l'anomalie (robuste,
marginale, normale), et l'indication de concordance inter-modèles. Un tableau
récapitulatif ordonne les $N$ périodes les plus anomales par score décroissant,
avec la mention des événements documentés concordants et de l'accord
inter-modèles pour chaque alerte.

% ════════════════════════════════════════════════════════════════
\section{Conclusion}

\subsection{Bilan du document}

Ce document a établi la méthodologie complète du Modèle~1 de l'étude
comparative, en trois contributions structurées.

La première contribution est la justification empirique du choix du LOF parmi
huit méthodes concurrentes. Le LOF est retenu comme meilleur compromis pour
le contexte BEAC ($N = 195$, $D = 55$, données financières à densité variable)
sur la base de sa supériorité empirique sur l'Isolation Forest sur données
financières de petit volume (F1 = 33.33\% vs 17.82\%, arXiv:2209.11686), de
son adéquation structurelle aux données multi-régimes, et de sa légitimité
documentée dans cinq publications institutionnelles BIS/BCE/Banca d'Italia.

La deuxième contribution est la documentation des cinq lacunes structurelles
du LOF avec leurs données empiriques et leurs implications pour le contexte
BEAC. L'écart moyen de $-33{,}98$ points de F1 entre le LOF et l'Anomaly
Transformer (Xu et al.\ 2022) quantifie le potentiel de gain de BiVAT sur
la lacune temporelle seule. Les quatre lacunes restantes (dimension, incertitude,
explicabilité, instabilité) définissent précisément les axes de contribution
de BiVAT documentés dans le document d'état de l'art.

La troisième contribution est le protocole de comparaison inter-modèles sans
labels ground truth, fondé sur quatre métriques d'évaluation non supervisée
(Precision@k, rang des anomalies connues, cohérence inter-modèles,
concordance systémique) adaptées au contexte de surveillance
macroprudentielle BEAC-CEMAC.

\subsection{Transition vers BiVAT}

Le Banco de Portugal (Costa et al.\ 2021) documente que les méthodes
traditionnelles comme le LOF ``will (in general) perform poorly'' au-delà
d'une certaine complexité des données de surveillance bancaire. Cette
observation institutionnelle, combinée aux cinq lacunes documentées en
section~5, constitue la motivation précise de BiVAT comme Modèle~2 : non
pas un remplacement arbitraire du LOF par une méthode plus complexe, mais
une réponse architecturale documentée, lacune par lacune, aux limites
structurelles d'un algorithme classique légitime dans son domaine de
compétence et clairement insuffisant au-delà.

% ════════════════════════════════════════════════════════════════
\section*{Bibliographie}
\addcontentsline{toc}{section}{Bibliographie}

\begin{enumerate}

\item Breunig, M.M., Kriegel, H.-P., Ng, R.T., \& Sander, J. (2000).
LOF: Identifying Density-Based Local Outliers. \textit{Proceedings of the
2000 ACM SIGMOD International Conference on Management of Data},
pp.~93--104. DOI:~10.1145/342009.335388.

\item Xu, J., Wu, H., Wang, J., \& Long, M. (2022). Anomaly Transformer:
Time Series Anomaly Detection with Association Discrepancy. \textit{ICLR
2022 (Spotlight)}.
\url{https://ise.thss.tsinghua.edu.cn/~mlong/doc/anomaly-transformer-iclr22.pdf}

\item Xu, Y. et al. (2023). MEMTO: Memory-guided Transformer for Multivariate
Time Series Anomaly Detection. arXiv:2312.02530.

\item Chandola, V., Banerjee, A., \& Kumar, V. (2009). Anomaly Detection:
A Survey. \textit{ACM Computing Surveys}, 41(3), Article~15.
DOI:~10.1145/1541880.1541882.

\item Zimek, A., Schubert, E., \& Kriegel, H.-P. (2012). A survey on
unsupervised outlier detection in high-dimensional numerical data.
\textit{Statistical Analysis and Data Mining}, 5(5), 363--387.
DOI:~10.1002/sam.11161.

\item Liu, F.T., Ting, K.M., \& Zhou, Z.-H. (2008). Isolation Forest.
\textit{Proceedings of the 8th IEEE International Conference on Data Mining
(ICDM)}, pp.~413--422.

\item Goldstein, M., \& Uchida, S. (2016). A Comparative Evaluation of
Unsupervised Anomaly Detection Algorithms for Multivariate Data. \textit{PLoS
ONE}, 11(4): e0152173. DOI:~10.1371/journal.pone.0152173.

\item Candès, E.J., Li, X., Ma, Y., \& Wright, J. (2011). Robust Principal
Component Analysis? \textit{Journal of the ACM}, 58(3), Article~11.
DOI:~10.1145/1970392.1970395.

\item Cleveland, R.B., Cleveland, W.S., McRae, J.E., \& Terpenning, I. (1990).
STL: A Seasonal-Trend Decomposition Procedure Based on Loess.
\textit{Journal of Official Statistics}, 6(1), 3--73.

\item Van Buuren, S., \& Groothuis-Oudshoorn, K. (2011). mice: Multivariate
Imputation by Chained Equations in R. \textit{Journal of Statistical Software},
45(3), 1--67. DOI:~10.18637/jss.v045.i03.

\item Leys, C., Ley, C., Klein, O., Bernard, P., \& Licata, L. (2013).
Detecting outliers: Do not use standard deviation around the mean, use absolute
deviation around the median. \textit{Journal of Experimental Social Psychology},
49(4), 764--766. DOI:~10.1016/j.jesp.2013.03.013.

\item Kriegel, H.-P., Kröger, P., Schubert, E., \& Zimek, A. (2009). LoOP:
Local Outlier Probabilities. \textit{Proceedings of CIKM 2009}, pp.~1649--1652.
DOI:~10.1145/1645953.1646195.

\item Kriegel, H.-P., Kröger, P., Schubert, E., \& Zimek, A. (2011).
Interpreting and Unifying Outlier Scores. \textit{Proceedings of SDM 2011},
pp.~13--24. DOI:~10.1137/1.9781611972818.2.

\item Tukey, J.W. (1977). \textit{Exploratory Data Analysis}.
Addison-Wesley, Reading (MA).

\item Schölkopf, B., Platt, J.C., Shawe-Taylor, J., Smola, A.J., \&
Williamson, R.C. (2001). Estimating the Support of a High-Dimensional
Distribution. \textit{Neural Computation}, 13(7), 1443--1471.

\item Tang, J., Chen, Z., Fu, A.W.-C., \& Cheung, D.W. (2002). Enhancing
Effectiveness of Outlier Detections for Low Density Patterns. \textit{PAKDD
2002}, pp.~535--548.

\item Papadimitriou, S., Kitagawa, H., Gibbons, P.B., \& Faloutsos, C. (2003).
LOCI: Fast Outlier Detection Using the Local Correlation Integral.
\textit{Proceedings of ICDE 2003}, pp.~315--326.

\item Jin, W., Tung, A.K.H., Han, J., \& Wang, W. (2006). Ranking Outliers
Using Symmetric Neighborhood Relationship. \textit{PAKDD 2006}, pp.~577--593.

\item Zhang, K., Hutter, M., \& Jin, H. (2009). A New Local Distance-Based
Outlier Detection Approach for Scattered Real-World Data. \textit{PAKDD 2009},
pp.~813--822.

\item Lazarevic, A., \& Kumar, V. (2005). Feature Bagging for Outlier
Detection. \textit{Proceedings of the 11th ACM SIGKDD}, pp.~157--166.

\item He, X. et al. (2024). VAEAT: Variational AutoEncoder with Adversarial
Training. \textit{Information Sciences}, 676, 120852.

\item Su, Y. et al. (2019). Robust Anomaly Detection for Multivariate Time
Series through Stochastic Recurrent Neural Network. \textit{ACM SIGKDD},
pp.~2828--2837.

\item Bezanson, S. et al. (2026). Comparative imputation methods for
clinical time series. \textit{Nature Scientific Reports}.
DOI:~10.1038/s41598-026-39035-z.

\item BIS / IFC (2024). A scalable, explainable machine learning approach
for anomaly detection in banking data. \textit{IFC Bulletin~66}.
\url{https://www.bis.org/ifc/publ/ifcb66\_08.pdf}

\item Maddaloni, A. et al. (2022). Anomaly Detection in Banking Supervisory
Data. Bank of Italy Occasional Paper~689. SSRN:~4109500.

\item Accornero, G., \& Boscariol, F. (2021). Machine learning for anomaly
detection in datasets with categorical variables. \textit{BIS IFC Bulletin~57},
Section~5.
\url{https://www.bis.org/ifc/publ/ifcb57\_05.pdf}

\item Costa, A., Fonseca, J., \& Maurício, A. (2021). Novel methodologies for
data quality management: Anomaly detection in the Portuguese Central Credit
Register. \textit{BIS IFC Bulletin~57}, Section~29.
\url{https://www.bis.org/ifc/publ/ifcb57\_29.pdf}

\item BIS / IFC (2022). Machine learning applications in central banking.
\textit{IFC Bulletin~57}.
\url{https://www.bis.org/ifc/publ/ifcb57\_01\_rh.pdf}

\item BIS (2024). Finding a needle in a haystack: a ML framework for anomaly
detection in payment systems. \textit{BIS Bulletin~84}.
\url{https://www.bis.org/publ/bisbull84.pdf}

\item BIS / FSI (2024). How regulators can address AI explainability.
\url{https://www.bis.org/fsi/fsipapers24.pdf}

\item ECB (2025). Guide to Internal Models (ML chapter).
\url{https://www.bankingsupervision.europa.eu}

\item Moukoko Ekambi, G.M. (2025). Détection de Variations Atypiques dans
les Séries Temporelles Financières par Architectures Hybrides VAE-Transformer :
État de l'art outillé et justification architecturale de BiVAT. Rapport de
stage, BEAC / ENSPD / Université de Douala, LIDIA Laboratory.

\end{enumerate}

\vspace{2cm}
\hrule
\vspace{0.5cm}
\noindent{\small\textit{Document rédigé dans le cadre du stage de recherche
à la Banque des États de l'Afrique Centrale (BEAC), M2 Data Science \& IA,
ENSPD / Université de Douala --- LIDIA Laboratory. Août 2026.}}

\end{document}